\documentclass[11pt,a4paper]{article}
\usepackage[utf8]{inputenc}
\renewcommand{\abstractname}{Sommario}
\renewcommand{\figurename}{Figura}
\AtBeginDocument{\renewcommand{\proofname}{Dimostrazione}}
\usepackage{amsmath,amssymb,amsthm}
\usepackage{graphicx}
\usepackage{booktabs}
\usepackage[margin=2.4cm]{geometry}
\usepackage{hyperref}

\newtheorem{prop}{Proposizione}
\theoremstyle{remark}
\newtheorem*{oss}{Osservazione}

\newcommand{\R}{\mathbb{R}}
\newcommand{\bb}{\mathbf{b}}
\newcommand{\x}{\mathbf{x}}
\newcommand{\z}{\mathbf{z}}

\title{Stili di gioco e gruppi di lavoro bilanciati\\
\large PCA, metodo delle potenze e decomposizione della varianza\\ applicati a una squadra di rugby}
\author{Matteo Guerrucci}
\date{Progetto di Analisi Numerica --- settembre 2026}

\begin{document}
\maketitle

\begin{abstract}
Una rosa di $N=25$ giocatori \`e descritta da $D=10$ statistiche per 80 minuti.
Con l'analisi delle componenti principali (PCA) si individuano pochi \emph{assi di stile di gioco};
le componenti sono calcolate con un metodo delle potenze implementato da zero, con deflazione,
e verificate contro la SVD. Nello spazio degli stili si formano poi 5 gruppi di lavoro
\emph{bilanciati} minimizzando la varianza tra i gruppi: per la decomposizione della varianza
questo equivale a mescolare in ogni gruppo giocatori forti e deboli su ogni asse.
I dati sono sintetici.
\end{abstract}

%=====================================================================
\section{Richiami: valore atteso, varianza, covarianza}
Riferimento: M.~P.~Deisenroth, A.~A.~Faisal, C.~S.~Ong, \emph{Mathematics for Machine Learning},
Cambridge University Press, 2020 (cap.~6, statistiche riassuntive; cap.~10, PCA). Il libro usa la
convenzione $X\in\R^{D\times N}$ (dati per colonne) e il fattore $1/N$; qui $X\in\R^{N\times D}$
(una riga per giocatore) e $1/(N-1)$, lo stimatore non distorto.

\paragraph{Valore atteso e varianza.}
Per una variabile aleatoria $x$ con media $\mu=\mathbb{E}[x]$:
\[
\mathbb{V}[x]=\mathbb{E}\big[(x-\mu)^2\big]=\mathbb{E}[x^2]-\big(\mathbb{E}[x]\big)^2,
\qquad \sigma=\sqrt{\mathbb{V}[x]}\ \ \text{(deviazione standard)}.
\]
Per un vettore aleatorio $\x\in\R^D$ con media $\boldsymbol\mu$, la matrice di covarianza \`e
\[
\Sigma=\operatorname{Cov}[\x]=\mathbb{E}\big[(\x-\boldsymbol\mu)(\x-\boldsymbol\mu)^T\big],
\qquad \Sigma_{ij}=\operatorname{Cov}[x_i,x_j],\quad \Sigma_{jj}=\mathbb{V}[x_j].
\]
\paragraph{Versioni empiriche.} Da $N$ osservazioni $\x_1,\dots,\x_N$:
\[
\bar\x=\frac1N\sum_n\x_n,\qquad
C=\frac{1}{N-1}\sum_n(\x_n-\bar\x)(\x_n-\bar\x)^T=\frac{1}{N-1}X^TX\ \ (X\text{ centrata}).
\]

\paragraph{Perch\'e si centra prima: un fatto numerico.}
Le due espressioni della varianza sono uguali in aritmetica esatta, ma non in floating point.
$\mathbb{E}[x^2]-(\mathbb{E}[x])^2$ sottrae due numeri grandi e quasi uguali: \`e una
\textbf{cancellazione catastrofica}. Con i valori $\{4,7,13,16\}$ (varianza esatta $22.5$) spostati
di $10^9$, la formula a una passata d\`a $-128$ (una varianza negativa), quella a due passate
$\mathbb{E}[(x-\mu)^2]$ d\`a $22.5$. Per questo il codice centra sempre i dati prima di calcolare $C$.

\paragraph{Trasformazioni lineari.}
Se $\mathbf y=A\x+\mathbf c$, allora
\[
\mathbb{E}[\mathbf y]=A\boldsymbol\mu+\mathbf c,\qquad \operatorname{Cov}[\mathbf y]=A\,\Sigma\,A^T .
\]
In particolare, per la proiezione su una direzione $\bb$ ($A=\bb^T$):
\[
\boxed{\ \mathbb{V}[\bb^T\x]=\bb^T\Sigma\,\bb\ }
\]
\`e il punto di partenza della PCA: la varianza lungo una direzione \`e una forma quadratica nella
matrice di covarianza. Con $B^T$ al posto di $\bb^T$: $\operatorname{Cov}[B^T\x]=B^T\Sigma B$.

\paragraph{Geometria delle variabili aleatorie.}
Per variabili a media nulla, $\langle x,y\rangle:=\operatorname{Cov}[x,y]$ \`e un prodotto scalare.
La norma indotta \`e la deviazione standard, $\|x\|=\sqrt{\operatorname{Cov}[x,x]}=\sigma_x$, e
\[
\cos\theta=\frac{\langle x,y\rangle}{\|x\|\,\|y\|}=\frac{\operatorname{Cov}[x,y]}{\sigma_x\sigma_y}=\rho(x,y),
\]
cio\`e \textbf{la correlazione \`e il coseno dell'angolo} tra le due variabili.
Variabili \textbf{scorrelate} sono \textbf{ortogonali}. In questa geometria la PCA cerca direzioni
ortogonali: gli score lungo componenti diverse sono scorrelati, perch\'e
$\operatorname{Cov}[\bb_i^T\x,\bb_j^T\x]=\bb_i^T\Sigma\bb_j=\lambda_j\bb_i^T\bb_j=0$ per $i\ne j$.
Inoltre $\operatorname{tr}(\Sigma)=\sum_j\mathbb{V}[x_j]$ \`e la varianza totale (la somma dei quadrati
delle ``lunghezze'' delle variabili).

%=====================================================================
\section{Dati e standardizzazione}
Sia $X_{\text{raw}}\in\R^{N\times D}$ la matrice dei dati (una riga per giocatore).
Ogni colonna viene centrata e divisa per la sua deviazione standard:
\[
X_{nj}=\frac{(X_{\text{raw}})_{nj}-\mu_j}{\sigma_j}.
\]
Senza standardizzazione le statistiche con valori grandi (metri guadagnati $\approx 50$)
dominerebbero quelle con valori piccoli (offload $\approx 1$) solo per una questione di unit\`a di misura.

\paragraph{Matrice di covarianza.}
\[
C=\frac{1}{N-1}X^{T}X\in\R^{D\times D}.
\]
Con dati standardizzati $C$ \`e la matrice di correlazione: $C_{jj}=1$ e $\operatorname{tr}(C)=D$.
$C$ \`e \textbf{simmetrica} e \textbf{semidefinita positiva}; per il teorema spettrale ha autovalori
reali $\lambda_1\ge\dots\ge\lambda_D\ge 0$ e autovettori ortonormali.

\paragraph{Lettura geometrica.}
Per colonne centrate $a,b$, la covarianza $\operatorname{cov}(a,b)=a^Tb/(N-1)$ \`e un prodotto scalare,
la deviazione standard \`e la norma associata e la correlazione
$\rho=\operatorname{cov}(a,b)/(\sigma_a\sigma_b)$ \`e il \textbf{coseno dell'angolo} tra le due colonne.
La nuvola dei giocatori in $\R^D$ ha la forma di un ellissoide: gli autovettori di $C$ ne sono gli assi,
gli autovalori la varianza lungo ciascun asse.

%=====================================================================
\section{PCA come problema di massima varianza}
Si cerca una base ortonormale $B=[\bb_1,\dots,\bb_M]\in\R^{D\times M}$, $B^TB=I$
($\bb_i^T\bb_j=\delta_{ij}$), e si proiettano i dati: $\z_n=B^T\x_n$.

\paragraph{Varianza lungo una direzione.}
Se $z_n=\bb^T\x_n$, allora
\[
\operatorname{Var}(z)=\frac{1}{N-1}\sum_n z_n^2=\frac{1}{N-1}(X\bb)^T(X\bb)=\bb^TC\bb .
\]

\subsection{Prima componente}
\[
\max_{\bb}\ \bb^TC\bb\qquad\text{con}\qquad \bb^T\bb=1 .
\]
Lagrangiana: $\ L(\bb,\lambda)=\bb^TC\bb+\lambda(1-\bb^T\bb)$.
Poich\'e $C$ \`e simmetrica, $\nabla_{\bb}(\bb^TC\bb)=(C+C^T)\bb=2C\bb$, quindi
\[
\frac{\partial L}{\partial\bb}=2C\bb-2\lambda\bb=0\ \Longrightarrow\ C\bb=\lambda\bb,
\qquad
\frac{\partial L}{\partial\lambda}=1-\bb^T\bb=0 .
\]
I punti stazionari sono gli autovettori (normalizzati) di $C$. Moltiplicando $C\bb=\lambda\bb$ per $\bb^T$:
\[
\bb^TC\bb=\lambda\,\bb^T\bb=\lambda,
\]
cio\`e \textbf{la varianza lungo un autovettore \`e il suo autovalore}.
Il massimo si ottiene con l'autovalore pi\`u grande $\lambda_1$: $\bb_1$ \`e la prima componente principale.

\subsection{Seconda componente}
\[
\max_{\bb_2}\ \bb_2^TC\bb_2\qquad\text{con}\qquad \bb_2^T\bb_2=1,\ \ \bb_2^T\bb_1=0 .
\]
\[
L=\bb_2^TC\bb_2+\lambda_2(1-\bb_2^T\bb_2)-\mu\,\bb_2^T\bb_1,
\qquad
\frac{\partial L}{\partial\bb_2}=2C\bb_2-2\lambda_2\bb_2-\mu\bb_1=0 .
\]
Moltiplicando a sinistra per $\bb_1^T$: il primo termine d\`a
$2\bb_1^TC\bb_2=2(C\bb_1)^T\bb_2=2\lambda_1\bb_1^T\bb_2=0$, il secondo $-2\lambda_2\bb_1^T\bb_2=0$,
il terzo $-\mu\,\bb_1^T\bb_1=-\mu$. Quindi $\mu=0$ e $C\bb_2=\lambda_2\bb_2$.
Poich\'e $\bb_2\perp\bb_1$, tra gli autovettori rimasti il massimo \`e $\lambda_2$ (supponendo $\lambda_1>\lambda_2$).
Per induzione, la $j$-esima componente \`e l'autovettore di $\lambda_j$.

\subsection{Varianza spiegata}
La varianza totale non dipende dalla base ortonormale:
\[
\sum_{j=1}^{D}\operatorname{Var}(X_{:,j})=\operatorname{tr}(C)=\sum_{j=1}^{D}\lambda_j .
\]
La frazione spiegata dalle prime $M$ componenti \`e $\sum_{j\le M}\lambda_j\big/\sum_{j}\lambda_j$.

\subsection{Legame con la SVD}
Se $X=U\Sigma V^T$, allora $X^TX=V\Sigma^2V^T$, quindi
\[
\lambda_j=\frac{\sigma_j^2}{N-1},\qquad \bb_j=\text{colonna $j$ di }V .
\]
In pratica si preferisce la SVD di $X$ al calcolo degli autovalori di $X^TX$ perch\'e
$K_2(X^TX)=K_2(X)^2$: formare $X^TX$ eleva al quadrato il numero di condizionamento.

\subsection{Seconda lettura: minimo errore di ricostruzione}
La PCA si pu\`o ottenere anche senza parlare di varianza (MML, par.~10.3). Proietto ogni $\x_n$ sul
sottospazio generato da $B$ e lo ricostruisco: $\tilde\x_n=BB^T\x_n$ ($BB^T$ \`e il proiettore
ortogonale). Cerco la base che rende minimo l'errore medio di ricostruzione
\[
J_M=\frac{1}{N-1}\sum_n\|\x_n-\tilde\x_n\|^2 .
\]
Per Pitagora, $\|\x_n\|^2=\|\tilde\x_n\|^2+\|\x_n-\tilde\x_n\|^2$, quindi
\[
\underbrace{\frac{1}{N-1}\sum_n\|\x_n\|^2}_{\operatorname{tr}(C)\ \text{(fissa)}}
=\underbrace{\sum_{j\le M}\bb_j^TC\bb_j}_{\text{varianza catturata}}+J_M .
\]
Minimizzare l'errore di ricostruzione equivale a massimizzare la varianza catturata: le due letture
danno la \textbf{stessa} base, e all'ottimo
\[
J_M=\sum_{j>M}\lambda_j ,
\]
la somma degli autovalori scartati. \`E la stessa affermazione di Eckart--Young (Sezione~\ref{sec:rango}), scritta
con gli autovalori di $C$ invece che con i valori singolari di $X$.

\paragraph{I passi della PCA in pratica} (MML, par.~10.6): (1) sottrarre la media; (2) dividere per
la deviazione standard; (3) calcolare autovalori e autovettori di $C$ (o la SVD di $X$); (4) proiettare,
$\z_n=B^T\x_n$; (5) per tornare nello spazio originale, $\tilde\x_n=B\z_n$, poi moltiplicare per
$\sigma_j$ e aggiungere $\mu_j$.

\begin{oss}
Se i giocatori fossero meno delle statistiche ($N<D$), $C$ avrebbe al pi\`u $N-1$ autovalori non
nulli, e conviene lavorare con la matrice $N\times N$ $XX^T$: se $XX^T\mathbf u=\mu\mathbf u$ allora
$X^TX(X^T\mathbf u)=\mu\,X^T\mathbf u$ (MML, par.~10.5). Con $N=25$, $D=10$ non serve.
\end{oss}

%=====================================================================
\section{Metodo delle potenze}
\paragraph{Algoritmo.}
\begin{center}
\begin{tabular}{l}
$\mathbf v\leftarrow\mathbf v_0/\|\mathbf v_0\|$ \quad ($\mathbf v_0$ casuale)\\
per $k=1,2,\dots$:\quad $\mathbf w=C\mathbf v$;\quad $\lambda^{(k)}=\mathbf v^T\mathbf w$;\quad
$\mathbf v\leftarrow\mathbf w/\|\mathbf w\|$\\
\quad stop se $|\lambda^{(k)}-\lambda^{(k-1)}|<\text{tol}\cdot|\lambda^{(k)}|$
\end{tabular}
\end{center}
Il quoziente di Rayleigh $\lambda^{(k)}=\mathbf v^TC\mathbf v$ (con $\|\mathbf v\|=1$) si ottiene
dallo stesso prodotto $\mathbf w=C\mathbf v$: un solo prodotto matrice--vettore per iterazione.
Il vettore iniziale casuale ha, con probabilit\`a 1, una componente non nulla lungo $\bb_1$.

\paragraph{Convergenza.}
Scriviamo $\mathbf v_0=\sum_i c_i\bb_i$ con $c_1\neq 0$. Allora
\[
C^k\mathbf v_0=\sum_i c_i\lambda_i^k\bb_i
=c_1\lambda_1^k\Big(\bb_1+\sum_{i\ge 2}\frac{c_i}{c_1}\Big(\frac{\lambda_i}{\lambda_1}\Big)^{k}\bb_i\Big).
\]
Se $|\lambda_2|<|\lambda_1|$ la direzione converge a $\bb_1$ con errore $O(|\lambda_2/\lambda_1|^k)$.
Poich\'e $C$ \`e simmetrica, l'errore sul quoziente di Rayleigh \`e $O(|\lambda_2/\lambda_1|^{2k})$.
La velocit\`a dipende dal \textbf{distacco} tra i primi due autovalori: se $\lambda_2\approx\lambda_1$
la convergenza \`e lenta; se $\lambda_2=\lambda_1$ non c'\`e una direzione privilegiata.

\paragraph{Deflazione.}
\[
C_1=C-\lambda_1\bb_1\bb_1^T .
\]
Per l'ortonormalit\`a degli autovettori, $C_1\bb_1=0$ e $C_1\bb_i=\lambda_i\bb_i$ per $i\ge 2$:
il metodo delle potenze su $C_1$ converge a $\bb_2$, con velocit\`a $|\lambda_3/\lambda_2|$.
La deflazione \`e la versione numerica del vincolo $\bb_2\perp\bb_1$ della lagrangiana.

\paragraph{Risultati.}
\begin{center}
\begin{tabular}{cccc}
\toprule
$j$ & $\lambda_j$ (potenze) & rapporto che comanda & iterazioni\\
\midrule
1 & 3.94498513 & $\lambda_2/\lambda_1=0.75$ & 58\\
2 & 2.97239586 & $\lambda_3/\lambda_2=0.35$ & 16\\
3 & 1.04652146 & $\lambda_4/\lambda_3=0.70$ & 39\\
\bottomrule
\end{tabular}
\end{center}
Gli autovalori coincidono con quelli della SVD fino all'ultima cifra mostrata.
Il fattore di riduzione dell'errore misurato per $\bb_1$ \`e $0.760$, contro il valore teorico
$\lambda_2/\lambda_1=0.754$ (Figura~\ref{fig:pot}). Le prime tre componenti spiegano il $79.6\%$ della varianza.

\begin{oss}
Il criterio d'arresto \`e sull'autovalore, che converge con velocit\`a doppia: quando $\lambda^{(k)}$ \`e
esatto a $10^{-12}$, l'autovettore lo \`e solo a circa $10^{-6}$. Se serve l'autovettore preciso
ci si ferma sul residuo $\|C\mathbf v-\lambda\mathbf v\|$.
\end{oss}

\begin{figure}[h]
\centering
\includegraphics[width=0.72\textwidth]{fig_potenze_convergenza.png}
\caption{Errore del metodo delle potenze per $\bb_1$: le curve sono parallele alle rette teoriche
$|\lambda_2/\lambda_1|^k$ (autovettore) e $|\lambda_2/\lambda_1|^{2k}$ (Rayleigh).}
\label{fig:pot}
\end{figure}

%=====================================================================
\section{Accelerare il metodo: shift e potenze inverse}
\paragraph{Shift.} $C-\sigma I$ ha gli stessi autovettori e autovalori $\lambda_i-\sigma$.
Per $\bb_1$ il rapporto di convergenza diventa $\max_{i\ge2}|\lambda_i-\sigma|/|\lambda_1-\sigma|$,
minimo per $\sigma^*=(\lambda_2+\lambda_D)/2$. Con i nostri dati $\sigma^*=1.504$: il rapporto scende
da $0.753$ a $0.601$ e le iterazioni da 47 a 29.

\paragraph{Potenze inverse con shift.} $(C-\sigma I)^{-1}$ ha autovalori $1/(\lambda_i-\sigma)$: domina
l'autovalore pi\`u vicino a $\sigma$, con rapporto
$|\lambda_j-\sigma|/|\lambda_{\text{vicino}}-\sigma|$. Non si calcola l'inversa: a ogni passo si risolve
$(C-\sigma I)\mathbf w=\mathbf v$. La fattorizzazione $LU$ (con pivoting) si calcola \textbf{una volta},
e ogni iterazione costa due sostituzioni triangolari, $O(D^2)$ invece di $O(D^3)$.
\begin{center}
\begin{tabular}{ccccc}
\toprule
cerco & $\sigma$ & rapporto & iterazioni & errore su $\lambda$\\
\midrule
$\lambda_1$ & 3.8 & 0.175 & 9 & $2\cdot10^{-14}$\\
$\lambda_2$ & 3.0 & 0.029 & 5 & $3\cdot10^{-15}$\\
$\lambda_3$ & 1.0 & 0.171 & 8 & $1\cdot10^{-14}$\\
$\lambda_{10}$ & 0 & 0.476 & 20 & $2\cdot10^{-15}$\\
\bottomrule
\end{tabular}
\end{center}
Con $\sigma=0$ si ottiene l'autovalore pi\`u piccolo (potenze inverse ``pure'').
Scegliendo $\sigma$ si pu\`o calcolare \emph{qualunque} autovalore, senza deflazione.

\paragraph{A cosa serve nel progetto: il condizionamento.}
Con una matrice $10\times10$ la velocit\`a non \`e un problema. L'uso concreto \`e un altro:
metodo delle potenze e potenze inverse con $\sigma=0$ danno gli autovalori estremi, e per una
matrice simmetrica definita positiva
\[
K_2(C)=\frac{\lambda_{\max}}{\lambda_{\min}}=\frac{3.945}{0.0364}=108.4,
\qquad
K_2(X)=\frac{\sigma_{\max}}{\sigma_{\min}}=10.41,\qquad K_2(X)^2=108.4 .
\]
\`E la verifica numerica di $K_2(X^TX)=K_2(X)^2$ e quantifica perch\'e si preferisce la SVD di $X$.
Inoltre $\lambda_{\min}\approx 0$ dice che lungo $\bb_{10}$ i giocatori quasi non variano: esiste una
combinazione di statistiche quasi costante, cio\`e statistiche quasi ridondanti.

%=====================================================================
\section{Affidabilit\`a degli assi}\label{sec:aff}
Si perturbano i dati con rumore gaussiano pari al $10\%$ della deviazione standard (300 prove) e si
misura l'angolo tra ogni autovettore originale e quello perturbato.
\begin{center}
\begin{tabular}{ccccc}
\toprule
asse & $\lambda_j$ & gap con il vicino & rotazione mediana & rotazione 90\%\\
\midrule
$\bb_1$ & 3.945 & 0.973 & $2.5^\circ$ & $5.1^\circ$\\
$\bb_2$ & 2.972 & 0.973 & $2.9^\circ$ & $5.3^\circ$\\
$\bb_3$ & 1.047 & 0.318 & $5.3^\circ$ & $9.2^\circ$\\
$\bb_4$ & 0.728 & 0.178 & $8.3^\circ$ & $13.6^\circ$\\
\bottomrule
\end{tabular}
\end{center}
Pi\`u il gap tra un autovalore e i suoi vicini \`e piccolo, pi\`u l'autovettore \`e \textbf{mal condizionato}:
una piccola perturbazione dei dati lo fa ruotare molto. Al limite $\lambda_j=\lambda_{j+1}$ l'autovettore non
\`e nemmeno definito (qualunque direzione del piano \`e un autovettore). \`E lo stesso gap che governa la
velocit\`a del metodo delle potenze: \emph{convergenza lenta e instabilit\`a hanno la stessa causa}.
Per questo si interpretano solo $\bb_1$ e $\bb_2$.
\begin{figure}[h]
\centering
\includegraphics[width=0.55\textwidth]{fig_sensibilita.png}
\caption{Rotazione degli autovettori con rumore del 10\% sui dati: cresce al diminuire del gap.}
\end{figure}

%=====================================================================
\section{Rango $k$ e dati mancanti}\label{sec:rango}
\paragraph{Eckart--Young.} Tra le matrici di rango $k$, la pi\`u vicina a $X$ \`e
$X_k=U_k\Sigma_kV_k^T$, con
\[
\|X-X_k\|_2=\sigma_{k+1},\qquad \|X-X_k\|_F=\Big(\sum_{i>k}\sigma_i^2\Big)^{1/2}.
\]
Verificato numericamente per ogni $k$. Poich\'e $\sum_i\sigma_i^2=\|X\|_F^2$, vale
$\big(\|X-X_k\|_F/\|X\|_F\big)^2=1-\text{varianza spiegata}$: con $k=3$ l'errore relativo \`e
$45\%=\sqrt{1-0.796}$.

\paragraph{Imputazione (i dati ``disordinati'').} Due giocatori hanno 4 statistiche mancanti. Invece di
sostituirle con la media, si usa la struttura di basso rango: le statistiche sono correlate, quindi quelle
osservate informano su quelle mancanti. Iterazione di punto fisso: si parte con i buchi uguali alla media;
a ogni passo si calcola l'approssimazione di rango $k$ e si sostituiscono \emph{solo} i valori mancanti
con quelli ricostruiti; ci si ferma quando i valori imputati non cambiano pi\`u.
Errore sui valori mancanti (RMSE, in deviazioni standard):
\begin{center}
\begin{tabular}{lccccc}
\toprule
metodo & media & $k=1$ & $k=2$ & $k=3$ & $k=4$\\
\midrule
RMSE & 1.19 & 0.95 & \textbf{0.75} & 0.96 & 3.16 (non converge)\\
\bottomrule
\end{tabular}
\end{center}
Con $k$ piccolo l'imputazione migliora la media del $37\%$. Con $k$ troppo grande il modello ha abbastanza
libert\`a da adattarsi ai dati osservati senza vincolare quelli mancanti: il punto fisso diventa lentissimo o
non converge e l'errore esplode. La convergenza, quando c'\`e, \`e lineare (fattore $0.915$ per $k=3$).
\begin{figure}[h]
\centering
\includegraphics[width=0.9\textwidth]{fig_eckart_young.png}\\[4pt]
\includegraphics[width=0.9\textwidth]{fig_imputazione.png}
\caption{Sopra: autovalori ed errore di rango $k$. Sotto: imputazione con SVD iterativa e convergenza del punto fisso.}
\end{figure}

%=====================================================================
\section{Interpretazione degli assi}
Il nome di un asse si legge dai pesi (loadings) pi\`u grandi, positivi e negativi.
\begin{itemize}
\item $\bb_1$ \textbf{contatto $\leftrightarrow$ gioco al largo}: pesi positivi su ruck, placcaggi, touche;
negativi su metri guadagnati, difensori battuti, calci. Piloni e seconde linee da una parte, ali ed estremo dall'altra.
\item $\bb_2$ \textbf{portatore di palla $\leftrightarrow$ chi porta poco palla}: pesi positivi su portate, offload,
difensori battuti. Dal lato opposto ci sono sia i registi sia i piloni: l'interpretazione iniziale
``regista contro portatore'' \`e stata corretta dai dati.
\end{itemize}
\begin{oss}
Gli assi descrivono \textbf{stili}, non livelli: uno score basso su $\bb_1$ per un'ala non indica una carenza.
Per parlare di margini di miglioramento bisogna confrontare giocatori dello stesso ruolo (Sezione~\ref{sec:dir}).
Inoltre la terza componente \`e meno affidabile: $\lambda_3=1.05$ e $\lambda_4=0.73$ sono vicini, e un piccolo
distacco tra autovalori rende l'autovettore sensibile al rumore dei dati.
\end{oss}

%=====================================================================
\section{Gruppi di lavoro bilanciati}
Siano $\z_n\in\R^{K}$ gli score ($K=3$), con media della squadra nulla. Si dividono i giocatori in
$G=5$ gruppi da 5; $\bar\z_g$ \`e la media del gruppo $g$ e $n_g$ il numero di giocatori.

\begin{prop}[Decomposizione della varianza]
\[
\underbrace{\sum_n\|\z_n\|^2}_{SS_{\text{tot}}}
=\underbrace{\sum_g n_g\|\bar\z_g\|^2}_{SS_{\text{tra}}}
+\underbrace{\sum_g\sum_{n\in g}\|\z_n-\bar\z_g\|^2}_{SS_{\text{dentro}}} .
\]
\end{prop}
\begin{proof}
Per ogni gruppo si scrive $\z_n=(\z_n-\bar\z_g)+\bar\z_g$ e si sviluppa il quadrato:
\[
\sum_{n\in g}\|\z_n\|^2=\sum_{n\in g}\|\z_n-\bar\z_g\|^2+n_g\|\bar\z_g\|^2
+2\,\bar\z_g^T\sum_{n\in g}(\z_n-\bar\z_g).
\]
L'ultima somma \`e nulla per definizione di media. Sommando sui gruppi si ottiene la tesi.
\end{proof}

\paragraph{Obiettivo.}
$SS_{\text{tot}}$ \`e fissata dai dati, quindi
\[
\min SS_{\text{tra}}\iff\max SS_{\text{dentro}} :
\]
gruppi con la stessa media della squadra sono, necessariamente, gruppi internamente vari,
in cui giocatori forti e deboli su ogni asse sono mescolati.
Gli assi non vanno ripesati: lo score lungo $\bb_j$ ha gi\`a varianza $\lambda_j$.

\paragraph{Algoritmo di scambio.}
Si parte da una divisione casuale; si provano tutti gli scambi di due giocatori di gruppi diversi e si
applica quello che diminuisce di pi\`u $SS_{\text{tra}}$; ci si ferma quando nessuno scambio migliora.
L'algoritmo trova un \textbf{minimo locale}: accetta solo scambi migliorativi e non pu\`o superare
una configurazione in cui serve prima peggiorare. Le divisioni possibili sono circa
$25!/((5!)^5\,5!)\approx 5\cdot 10^{12}$, troppe da esplorare; si usano 30 partenze casuali e si tiene la migliore.

\paragraph{Risultati.}
$SS_{\text{tot}}=191.13$, $SS_{\text{tra}}=0.30$ ($0.16\%$), $SS_{\text{dentro}}=190.83$; la somma verifica
l'identit\`a a $10^{-14}$. Su 5000 divisioni casuali $SS_{\text{tra}}$ ha media $31.7$ e minimo $4.2$
(Figura~\ref{fig:gruppi}). $SS_{\text{tra}}$ non arriva a 0 perch\'e ogni gruppo dovrebbe avere media
esattamente nulla su tutti e tre gli assi con esattamente 5 giocatori: in generale non \`e possibile.

\begin{figure}[h]
\centering
\includegraphics[width=0.49\textwidth]{fig_gruppi.png}\hfill
\includegraphics[width=0.49\textwidth]{fig_confronto_casuali.png}
\caption{Sinistra: giocatori nel piano $(\bb_1,\bb_2)$ colorati per gruppo; le X sono le medie dei gruppi,
tutte vicine all'origine. Destra: $SS_{\text{tra}}$ di 5000 divisioni casuali contro la divisione ottimizzata.}
\label{fig:gruppi}
\end{figure}

%=====================================================================
\section{Direzioni di lavoro}\label{sec:dir}
Per ogni giocatore si confronta lo score con la media del suo ruolo. Su un asse $j$ il verso tipico
del ruolo \`e $s=\operatorname{segno}(\bar z_{\text{ruolo},j})$; l'asse si usa solo se il ruolo sta
chiaramente da una parte ($|\bar z_{\text{ruolo},j}|\ge 0.5\sqrt{\lambda_j}$). Il deficit \`e
\[
d_j=s\,\frac{z_{n,j}-\bar z_{\text{ruolo},j}}{\sqrt{\lambda_j}} ,
\]
e si segnala l'asse con deficit pi\`u negativo (sotto $-0.25$). Le statistiche da allenare sono quelle
con peso grande nel verso $s$ di $\bb_j$ che il ruolo usa davvero e in cui il giocatore \`e sotto la media
del ruolo; il mentore \`e il compagno di gruppo pi\`u forte su quelle statistiche.
Esempio: Seconda linea 1, deficit $-0.46$ su $\bb_1$, allenare ruck e placcaggi, mentore Terza linea 2.

%=====================================================================
\section{Limiti}
\begin{itemize}
\item Dati sintetici, $N=25$: le componenti oltre la seconda sono poco affidabili (Sezione~\ref{sec:aff}).
\item L'algoritmo di scambio \`e un'euristica: nessuna garanzia di ottimo globale.
\item La PCA descrive stili, non livelli: la direzione di lavoro richiede ipotesi aggiuntive
(confronto con il ruolo).
\item Sviluppi: dati reali di un campionato (assi stimati su molti giocatori, gruppi dentro una squadra);
per $N$ grande il prodotto $C\mathbf v=X^T(X\mathbf v)/(N-1)$ evita di formare $C$.
\end{itemize}

\begin{thebibliography}{9}
\bibitem{mml} M. P. Deisenroth, A. A. Faisal, C. S. Ong, \emph{Mathematics for Machine Learning}, Cambridge University Press, 2020. Cap. 6 (statistiche riassuntive, geometria delle variabili aleatorie) e cap. 10 (PCA).
\end{thebibliography}

\end{document}
